\subsection{Rationalizing Denominators with Multiple Terms} 
If there is more than one term in the denominator, it can be more complicated to rationalize the denominator.
\begin{Example}
Consider the fraction $\dfrac{7}{3-\sqrt{5}}$. If we multiply the numerator and denominator by $\sqrt{5}$, does this rationalize the denominator?
\begin{lpic}{}{2}
\fractextleft{0.25}{-1}{$\dfrac{7}{3-\sqrt{5}}={}$};
\end{lpic}
\makeblankfill{This doesn't work because}
\end{Example}
\noi To rationalize such denominators, it is valuable to recall the difference of squares special form.
\begin{fact}[Difference of Squares]
$$(x-y)(x+y)=x^2-y^2$$
\end{fact}
This gives us a tool we can use to rationalize denominators with two terms, one or both of which is a square root.
\begin{definition}[Conjugates]
For positive real numbers $x$ and $y$, the expressions
$$\sqrt{x}-\sqrt{y}
\quad\mbox{and}\quad
\sqrt{x}+\sqrt{y}$$
are called \term{conjugates}. This is useful because
$$\left(\sqrt{x}-\sqrt{y}\right)
\left(\sqrt{x}+\sqrt{y}\right)
=
\left(\sqrt{x}\right)^2
-
\left(\sqrt{y}\right)^2
=
x-y$$
In words, multiplying by the conjugate allows us to eliminate the square roots.
\end{definition}
\begin{Note}
This definition makes sense even when one term doesn't contain a radical, since any positive real number is the square root of its square.
\end{Note}
\begin{Example}
Identify the conjugate of $3-\sqrt{5}$. Then, simplify $\dfrac{7}{3-\sqrt{5}}$ by rationalizing the denominator.
\begin{lpic}{}{3}
\regtextleft{0.25}{-1}{The conjugate of $3-\sqrt{5}$ is};
\fractextleft{0.25}{-2}{$\dfrac{7}{3-\sqrt{5}}={}$};
\end{lpic}
\end{Example}
\begin{Note}
We can make similar moves for higher-order roots using other identities, like the sum or difference of cubes for cube roots.
\end{Note}